% Part: first-order-logic
% Chapter: syntax-and-semantics
% Section: Structures

\documentclass[../../../include/open-logic-section]{subfiles}

\begin{document}

\olfileid{fol}{syn}{str}

\olsection{پہلے درجے دیاں زباناں لئی \printtoken{P}{structure}}

\begin{explain}
پہلے درجے دیاں زباناں اپنے آپ وچ \emph{بنا تعبیر دے} ہوندیاں نیں:
!!{constant}s، !!{function}s تے !!{predicate}s نال کوئی خاص
معنی پہلے توں جڑے نہیں ہوندے۔ معنی \article{structure}
\emph{!!{structure}} متعین کرکے دتے جاندے نیں۔ اوہ
\emph{دائرہ} متعین کردی اے، یعنی اوہ اشیاء جنہاں ول
!!{constant}s اشارہ کردے نیں، جنہاں اُتے !!{function}s کم کردے
نیں، تے جنہاں اُتے مسوّرے چلدے نیں۔ نال ایہ وی متعین ہوندا اے
کہ کیہڑے !!{constant}s کیہڑیاں اشیاء ول اشارہ کردے نیں،
!!a{function} اشیاء نوں دوجیاں اشیاء ول کیویں بھیجدا اے، تے
!!{predicate}s کیہڑیاں اشیاء اُتے لاگو ہوندے نیں۔
!!^{structure}s منطق وچ \emph{معنوی} تصورات، مثلاً منطقی نتیجے،
اعتبار تے پورا ہون جوگتا دی بنیاد نیں۔ علمی عبارت وچ انہاں نوں
کدے ``ساختاں،'' کدے ``تعبیراں،'' تے کدے ``نمونے'' آکھیا جاندا اے۔
\end{explain}

\begin{defn}[!!^{structure}s]
\Article{structure} \emph{!!{structure}}~$\Struct M$، پہلے درجے دی
زبان $\Lang{L}$ لئی، ہیٹھ لکھیاں چیزاں توں بندی اے:
\begin{enumerate}
\item \emph{دائرہ:} اک غیر خالی سیٹ، $\Domain M$
\item \emph{!!{constant}s دی تعبیر:} ہر !!{constant}~$c$ لئی
  جیہڑا $\Lang{L}$ وچ اے، !!a{element} $\Assign{c}{M} \in \Domain M$
\item \emph{!!{predicate}s دی تعبیر:} ہر $n$-مقامی
  !!{predicate}~$R$ لئی جیہڑا $\Lang{L}$ وچ اے ($\eq$ توں علاوہ)، اک $n$-مقامی
  تعلق $\Assign{R}{M} \subseteq \Domain{M}^n$
\item \emph{!!{function}s دی تعبیر:} ہر $n$-مقامی
  !!{function}~$f$ لئی جیہڑا $\Lang{L}$ وچ اے، اک $n$-مقامی فنکشن $\Assign{f}{M}
  \colon \Domain{M}^n \to \Domain{M}$
\end{enumerate}
\end{defn}

\begin{ex}
حساب دی زبان لئی !!^a{structure}~$\Struct M$ وچ اک سیٹ،
$\Domain M$ دا اک عنصر $\Assign{\Obj 0}{M}$ جیہڑا
!!{constant}~$\Obj 0$ دی تعبیر اے، اک یک-مقامی فنکشن
$\Assign{\Obj \prime}{M} \colon \Domain{M} \to \Domain M$،
دو دو-مقامی فنکشن $\Assign{\Obj +}{M}$ تے $\Assign{\Obj
    \times}{M}$ (دونیں $\Domain M^2 \to \Domain M$)، تے اک
دو-مقامی تعلق $\Assign{\Obj <}{M} \subseteq \Domain{M}^2$ ہوندے نیں۔

ایہو جہی ساخت دی اک صاف مثال ایہ اے:
\begin{enumerate}
\item $\Domain N = \Nat$
\item $\Assign{\Obj 0}{N} = 0$
\item $\Assign{\Obj \prime}{N}(n) = n + 1$ ہر $n \in \Nat$ لئی
\item $\Assign{\Obj +}{N}(n, m) = n + m$ ہر $n, m \in \Nat$ لئی
\item $\Assign{\Obj \times}{N}(n, m) = n\cdot m$ ہر $n, m \in \Nat$ لئی
\item $\Assign{\Obj <}{N} = \Setabs{\tuple{n, m}}{n \in \Nat, m \in
  \Nat, n < m}$
\end{enumerate}
ایس طرح تعریف کیتی ساخت~$\Struct N$ نوں، $\Lang L_A$ لئی،
\emph{حساب دا معیاری نمونہ} آکھدے آں، کیوں جے ایہ~$\Lang L_A$
دیاں غیر منطقی ثابتاں دی تعبیر ٹھیک اوہو جیہی کردی اے جیہڑی
تُسی توقع کردے او۔

پر~$\Lang L_A$ لئی ہور بہت ساریاں !!{structure}s وی ممکن نیں۔
مثلاً اسیں دائرے لئی صحیح عدداں دا سیٹ~$\Int$ لے سکدے آں،
قدرتی اعداد~$\Nat$ دی تھاں، تے $\Obj
0$، $\Obj \prime$، $\Obj +$، $\Obj \times$ تے $\Obj <$
دیاں تعبیراں اوہدے مطابق متعین کر سکدے آں۔ بلکہ~$\Lang L_A$
لئی ایسی ساختاں وی بنا سکدے آں جنہاں دا عدداں نال دور دا وی
واسطہ نہ ہووے۔
\end{ex}

\begin{ex}
ساخت~$\Struct M$ نوں، سیٹ نظریے دی زبان~$\Lang L_Z$ لئی،
بس اک سیٹ تے اک دو-مقامی تعلق چاہیدا اے۔ ایس لئی، خالص فنی
لحاظ نال، مثلاً لوکاں دا سیٹ تے تعلق ``$x$، $y$ توں وڈی عمر
دا اے''، $\Lang L_Z$ لئی !!a{structure} بن سکدا اے؛ ایویں
$\Nat$ تے $n \ge m$ دا تعلق وی، جدوں $n, m \in \Nat$ ہون۔

$\Lang L_Z$ لئی اک خاص دلچسپ !!{structure} اوہ اے جیہدے دائرے دے
!!{element}s خود سیٹ نیں تے $\Obj \in$ دی تعبیر واقعی تعلق
``$x$،~$y$ دا !!a{element} اے'' اے۔ ایہ \emph{وراثتاً محدود سیٹاں}
دی !!{structure}~$\Struct{{HF}}$ اے:
\begin{enumerate}
\item $\Domain{{{HF}}} = \emptyset \cup \Pow{\emptyset} \cup
  \Pow{\Pow{\emptyset}} \cup \Pow{\Pow{\Pow{\emptyset}}} \cup \dots$;
\item $\Assign{\Obj \in}{{{HF}}} = \Setabs{\tuple{x, y}}{x, y \in
  \Domain{{{HF}}}, x \in y}$.
\end{enumerate}
\end{ex}

\begin{digress}
کیہڑی شے نوں !!a{structure} منن دے ضابطے ساڈی منطق اُتے اثر
پاندے نیں۔ مثال لئی، خالی دائریاں دی ممانعت، مسوّر والے جملیاں
دے پورا اترن (یا سچ ہون) دی عام تعبیر ہِٹھ، ایہ یقینی بناندی اے
کہ $\lexists[x][(!A(x) \lor \lnot !A(x))]$ معتبر اے---یعنی
منطقی سچائی اے۔ تے ایہ شرط کہ سارے !!{constant}s دائرے دی
کسے شے ول اشارہ کرن، وجودی تعمیم نوں استدلال دا درست نمونہ
بناندی اے: $!A(a)$، ایس لئی $\lexists[x][!A(x)]$۔ جے اسیں
ناواں نوں دائرے توں باہر دیاں اشیاء ول اشارہ کرن دی اجازت
دیندے، یا اشارہ نہ کرن دی، تاں اسیں \emph{آزاد منطق} ول جا
رہے ہوندے؛ اوتھے وجودی تعمیم لئی اک وادھو مقدمہ چاہیدا اے:
$!A(a)$ تے $\lexists[x][\eq[x][a]]$، ایس لئی
$\lexists[x][!A(x)]$۔
\end{digress}

\end{document}
